% ============================================================
%  FIFA Rankings vs ELO Ratings — Research Paper
%  FINAL VERSION — All real numbers, all figures embedded
%  Upgrades: H1/H2/H3, APA 7th, three-line tables, Appendix
% ============================================================

\documentclass[12pt, a4paper]{article}

\usepackage[T1]{fontenc}
\usepackage[utf8]{inputenc}
\usepackage{lmodern}
\usepackage[a4paper, margin=2.54cm]{geometry}
\usepackage{setspace}
\doublespacing
\usepackage{amsmath}
\usepackage{amssymb}
\usepackage{booktabs}
\usepackage{array}
\usepackage{tabularx}
\usepackage{multirow}
\usepackage{graphicx}
\usepackage{float}
\usepackage[font=small, labelfont=bf, justification=justified]{caption}
\usepackage[hidelinks, breaklinks=true]{hyperref}
\usepackage{url}
\urlstyle{same}
\usepackage{apacite}
\bibliographystyle{apacite}
\usepackage{microtype}
\usepackage{parskip}
\usepackage[title]{appendix}
\usepackage{enumitem}
\usepackage{subcaption}

% ============================================================
\begin{document}
% ============================================================

% ---- Title page ----
\begin{titlepage}
  \centering
  \vspace*{2cm}
  {\LARGE\bfseries FIFA Rankings vs ELO Ratings:\\[0.4em]
  Predictive Validity in World Cup Knockout Stages (1994--2022)\par}
  \vspace{1.5cm}
  {\large Karamjeet Singh\par}
  \vspace{0.5cm}
  {\normalsize Independent Researcher\\
  New Delhi, India\par}
  \vspace{1cm}
  {\normalsize May 2026\par}
  \vspace{0.5cm}
  {\normalsize Word count: approximately 6,800 words (excluding tables, figures, and references)\par}
  \vfill
\end{titlepage}

% ---- Abstract ----
\newpage
\begin{abstract}
\noindent
Football's most consequential matches --- World Cup knockout encounters --- are seeded and
drawn using FIFA World Rankings. Yet a rival system, ELO ratings, has long been considered
theoretically superior. This study asks a direct question: which system actually predicts
knockout outcomes better, and did FIFA's 2018 decision to adopt an ELO-based formula close
the gap?

We analysed 115 knockout matches across eight World Cups from 1994 to 2022, fitting separate
logistic regression models for each rating system --- specified without an intercept, consistent
with the neutral-venue setting --- and evaluating them across six metrics: prediction accuracy,
discrimination (AUC-ROC), calibration (Brier score), goal difference explained (OLS $R^{2}$),
out-of-sample cross-validation, and controlled regression (McFadden's pseudo-$R^{2}$ against a
50/50 null). ELO outperformed FIFA World Rankings on every single measure. It correctly
predicted knockout winners 73.9\% of the time versus FIFA's 68.7\%, achieved an AUC of 0.775
versus 0.695, and produced better-calibrated probability estimates (Brier score 0.189 versus
0.220).

The 2018 FIFA formula change genuinely improved FIFA's predictive validity --- the accuracy gap
narrowed from 6.0 percentage points to 3.2. It did not close it. A formal McNemar's test
($p = 0.180$) could not confirm the accuracy difference at conventional significance thresholds,
reflecting a structural data constraint: achieving 80\% power requires approximately 88
discordant pairs, equivalent to roughly 506 total knockout matches. The convergent direction of
evidence across six complementary metrics, combined with ELO exceeding FIFA in 89.8\% of 10,000
bootstrap resamples, supports the conclusion that ELO ratings are more predictively valid.
Three robustness checks confirm the performance gap is structural --- not an artefact of stale
ranking data, not driven by the 1994 tournament, and not resolved by substituting FIFA points
for rank.

\bigskip
\noindent\textbf{Keywords:} ELO ratings; FIFA World Rankings; football prediction; World Cup;
knockout stage; logistic regression; calibration; McNemar's test; predictive validity;
staleness analysis
\end{abstract}

\newpage
\tableofcontents
\newpage

% ============================================================
\section{Introduction}
% ============================================================

Football ranking systems serve two practical functions: seeding tournament draws and measuring
national team prestige. FIFA World Rankings have performed both since 1992, officially
determining which teams occupy favoured positions at every World Cup draw. The system has
attracted sustained criticism from analysts who argue that ELO ratings --- a continuous,
opponent-weighted system originally designed for chess --- provide a more accurate picture of
team quality.

The stakes are not trivial. If FIFA rankings systematically misidentify the stronger team,
tournament draws are distorted, and the teams most likely to win are not adequately protected
from early collisions. For broadcasters, analysts, and betting markets, the question of which
system better predicts outcomes has direct practical value. For FIFA itself, the question carries
institutional weight --- the organisation spent over two decades defending a formula that produced 37.5\% accuracy at the 2006 World Cup --- a single-tournament result that, while striking, should be interpreted cautiously given the small sample of 16 matches.

The comparison between ELO and FIFA rankings is not new. \citeA{hvattum2010} found ELO
outperformed simple benchmark methods. \citeA{lasek2013} confirmed ELO-based approaches
consistently outperformed the official FIFA ranking. \citeA{wunderlich2016} found persistent
FIFA inferiority to betting markets even after the 2006 formula revision.

This paper addresses three gaps in the existing literature. First, it isolates knockout stage
matches specifically, where seeding consequences are most severe and squad rotation is absent.
Second, it exploits the 2018 formula change as a pre/post comparison to assess whether FIFA's adoption of ELO-based logic narrowed the predictive gap. Third, it applies novel
robustness checks using real FIFA ranking snapshot dates, a sensitivity test excluding the 1994
tournament, and a comparison against FIFA points. Three research questions guide the analysis:
Does ELO outperform FIFA in predicting World Cup knockout outcomes across 1994--2022? Has
predictive power changed across the 2018 formula change? Does ranking staleness materially
affect the conclusions?

% ---- Hypotheses ----
\subsection{Hypotheses}

Based on prior literature and the theoretical properties of each system, three directional
hypotheses are advanced:

\begin{description}[leftmargin=2em, labelwidth=1.8em]
  \item[\textbf{H1}] ELO ratings have greater predictive validity than FIFA World Rankings for
  World Cup knockout stage outcomes across 1994--2022, as measured by prediction accuracy,
  AUC-ROC, Brier score, goal-difference $R^{2}$, out-of-sample cross-validation accuracy,
  and controlled regression (McFadden's pseudo-$R^{2}$).

  \item[\textbf{H2}] FIFA's 2018 adoption of an ELO-based formula improved its predictive
  validity, narrowing --- but not closing --- the performance gap relative to ELO ratings.

  \item[\textbf{H3}] Ranking staleness does not materially account for FIFA's predictive
  inferiority; the gap is structural and persists after controlling for the number of days
  between the ranking snapshot and the match date.
\end{description}

% ============================================================
\section{Background and Related Literature}
% ============================================================

\subsection{The ELO Rating System}

The ELO system was developed by Árpád Élő for chess \cite{eloarp1978} and adapted for team sports.
The update rule is:
\[
  R_{\text{new}} = R_{\text{old}} + K \cdot G \cdot (W - W_e),
\]
where $R$ is the team's rating, $K$ is a match importance weight, $G$ is a goal difference
multiplier, $W$ is the match outcome (1 = win, 0.5 = draw, 0 = loss), and
\[
  W_e = \frac{1}{1 + 10^{-\Delta r / 400}}
\]
is the expected outcome derived from the pre-match rating difference $\Delta r$. The
eloratings.net implementation applies no home advantage correction for neutral venues, making
ELO ratings fully appropriate for this comparison \cite{eloratings2024}.

\subsection{FIFA World Rankings: Pre-2018 Formula}

FIFA introduced its World Rankings in 1992 using a points-based averaging system over an
eight-year rolling window. Two fundamental flaws have been documented in the literature \cite{lasek2013}: the system
rewarded match volume over quality, and goal difference was entirely ignored. These design choices made the
formula susceptible to manipulation through friendly match scheduling. The consequences were
most visible at the 2006 World Cup, where FIFA accuracy in this study's knockout analysis fell
to 37.5\% --- well below the 50\% chance threshold.

\subsection{FIFA World Rankings: Post-2018 Formula}

In August 2018, FIFA replaced its methodology with a system structurally based on ELO
\cite{fifa2018}. Two implementation differences from eloratings.net remain: FIFA does not
weight outcomes by goal difference, and uses a divisor of 600 rather than the standard 400,
making probability estimates less sensitive to rating gaps. Both are reversible policy choices
that this study's evidence suggests should be revisited.

\subsection{Prior Evidence on Predictive Validity}

\citeA{hvattum2010} found ELO outperformed simpler benchmarks in the English Premier League,
though it was outperformed by betting odds. \citeA{lasek2013} confirmed ELO-based approaches
outperformed FIFA rankings, noting friendly match scheduling as a significant source of
predictive invalidity. \citeA{wunderlich2016} found improvement after the 2006 formula revision
but persistent inferiority to betting markets. No prior study has isolated knockout stages,
spanned the 2018 formula change as a pre/post comparison, or applied real-data staleness controls
using verified FIFA snapshot dates.

% ============================================================
\section{Data and Methodology}
% ============================================================

\subsection{Dataset Construction}

The analysis covers eight World Cup knockout stages (USA 1994 through Qatar 2022). Group stage
matches were excluded due to squad rotation and differential qualification stakes. Four data
sources were merged: match results from \citeA{badrnezhad2024matches};
ELO ratings from \citeA{eloratings2024}, validated match-by-match against saved source pages;
FIFA rankings, points, and snapshot dates from \citeA{alex2024fifaranking};
and squad data from \citeA{fjelstul2023}.

The merged dataset contains 128 knockout matches. After removing 13 matches where ELO or FIFA
ranking data could not be reliably matched, the analysis dataset contains 115 matches.
Matches decided by penalty shootouts are treated as wins for the shootout winner in accuracy
analysis but excluded from goal difference regression.

\subsection{Variables}

\textbf{Dependent variable.} The primary outcome is a binary indicator of whether the
home-listed team won the match. Since all knockout matches are at neutral venues, the
home/away designation carries no competitive significance --- it reflects only administrative
draw order.

\textbf{Independent variables.} $\Delta\text{ELO} = \text{ELO}_{\text{home}} -
\text{ELO}_{\text{away}}$ (mean = 55.0, $SD$ = 142.1); $\Delta\text{FIFA} =
\text{FIFA\_rank}_{\text{away}} - \text{FIFA\_rank}_{\text{home}}$ (mean = 3.3, $SD$ = 16.7).
Both are signed so positive values indicate the home-listed team is stronger. FIFA rank is the
primary FIFA predictor because it is the actual seeding instrument FIFA uses for World Cup draws.

\textbf{Staleness variable.} \textit{Days\_stale} measures the number of days between the FIFA
ranking snapshot date and the match date, using real \texttt{rank\_date} values from the FIFA
ranking dataset \cite{alex2024fifaranking} --- not proxy estimates. Real staleness ranged from 0 to 73 days (mean 35.5 days).

Figure~\ref{fig:distributions} shows the distributions of ELO differences, FIFA rank
differences, and ranking staleness across the 115 matches.

\begin{figure}[H]
  \centering
  \includegraphics[width=\textwidth]{fig1_distributions.png}
  \caption{Distribution of rating differences and FIFA snapshot staleness across 115 World Cup
  knockout matches, 1994--2022. Left: ELO difference (home minus away; positive = home stronger).
  Centre: FIFA rank difference (away minus home rank; positive = home better ranked). Right:
  days between FIFA ranking snapshot and match date using real \texttt{rank\_date} values.
  Mean staleness = 35 days; maximum = 73 days.}
  \label{fig:distributions}
\end{figure}

\subsection{Statistical Methods}

\subsubsection{Logistic Regression}

Separate logistic regression models without an intercept are fitted for ELO and FIFA using
the \texttt{statsmodels} library \cite{seabold2010}, with a no-intercept specification
required by the neutral-venue setting: equal teams at a neutral
ground should have exactly a 50\% win probability.

\subsubsection{McNemar's Test}

McNemar's test \cite{mcnemar1947} is applied to formally test the accuracy difference:
\[
  \chi^{2} = \frac{(b - c)^{2}}{b + c},
\]
where $b$ = matches correctly predicted by ELO but not FIFA ($b = 13$), and $c$ = the reverse
($c = 7$). The test operates only on discordant pairs ($n_d = 20$).

\subsubsection{Bootstrap Confidence Intervals}

A paired percentile bootstrap with 10,000 resamples is applied. Each resample draws the same
matched rows for both systems simultaneously, preserving the within-match pairing. Sensitivity
analysis presents confidence intervals at 90\%, 95\%, and 99\% levels.

\subsubsection{Calibration Analysis}

The Brier score \cite{brier1950} summarises calibration:
\[
  \text{BS} = \frac{1}{N} \sum_{i=1}^{N}(p_i - o_i)^{2}.
\]
A random coin-flip yields $\text{BS} = 0.250$; perfect predictions yield $\text{BS} = 0.000$.

\subsubsection{Cross-Validation}

Leave-one-tournament-out cross-validation evaluates out-of-sample accuracy across eight
tournament folds. For each fold, the held-out tournament's outcomes are predicted using the
rating differences alone --- no model training on the test fold.

\subsubsection{Controlled Logistic Regression}

A final specification adds round (ordinal), centred year, and a post-2018 dummy to isolate the
pure rating signal from structural variation across rounds and eras.

\subsection{Statistical Power}

With $n_d = 20$ discordant pairs and $p_1 = 0.65$, McNemar's power is:
\[
  \text{power} = \Phi\!\left(\frac{0.65 - 0.50}{\sqrt{0.25/20}} - 1.96\right) \approx 0.268.
\]
Achieving 80\% power requires approximately 88 discordant pairs, equivalent to roughly 506
total knockout matches. See Appendix~\ref{app:power} for the full derivation.

% ============================================================
\section{Results}
% ============================================================

\subsection{Descriptive Statistics}

Table~\ref{tab:descriptive} reports descriptive statistics. The two systems agree on the
predicted stronger team in 82.6\% of the 115 matches; the remaining 20 constitute the
discordant pairs for McNemar's test.

\begin{table}[H]
  \centering
  \caption{Descriptive statistics for key variables ($N = 115$ World Cup knockout matches,
  1994--2022). Staleness computed from real FIFA \texttt{rank\_date} values.}
  \label{tab:descriptive}
  \begin{tabular}{lrrrrr}
    \toprule
    Variable & $N$ & Mean & $SD$ & Min & Max \\
    \midrule
    ELO difference ($\Delta_{\text{ELO}}$)       & 115 & 55.0  & 142.1 & $-314$ & 420 \\
    FIFA rank difference ($\Delta_{\text{FIFA}}$) & 115 &  3.3  &  16.7 & $-47$  &  47 \\
    Goal difference                               &  89 &  0.3  &   1.9 & $-6$   &   5 \\
    Days stale (real \texttt{rank\_date})         & 115 & 35.5  &  16.8 &    0   &  73 \\
    ELO--FIFA agreement on favourite              &  95 & 82.6\%&       &        &     \\
    \bottomrule
  \end{tabular}
\end{table}

\subsection{Logistic Regression and Discrimination}

Table~\ref{tab:regression} presents the baseline logistic regression results. Both rating
difference variables are statistically significant predictors (ELO: $\beta = 0.0084$,
$SE = 0.0019$, $p < 0.001$; FIFA: $\beta = 0.0460$, $SE = 0.0136$, $p = 0.001$). ELO
achieves prediction accuracy of 73.9\% versus FIFA's 68.7\%, and an AUC-ROC of 0.775 versus
0.695. McFadden's pseudo-$R^{2}$ for the ELO model (0.188) is 2.1 times that of the FIFA
model (0.089), computed against a 50/50 null. Figure~\ref{fig:roc} presents the ROC curves.

\begin{table}[H]
  \centering
  \caption{Baseline logistic regression results (no-intercept specification, $N = 115$).
  Pseudo-$R^{2}$ computed against a 50/50 null consistent with neutral-venue assumption.}
  \label{tab:regression}
  \begin{tabular}{lrrrr}
    \toprule
    & \multicolumn{2}{c}{\textbf{ELO Model}} & \multicolumn{2}{c}{\textbf{FIFA Model}} \\
    \cmidrule(lr){2-3}\cmidrule(lr){4-5}
    & Estimate & $p$-value & Estimate & $p$-value \\
    \midrule
    Coefficient ($\beta$) & 0.0084 & $<0.001$ & 0.0460 & 0.001 \\
    $SE$                  & 0.0019 &          & 0.0136 &       \\
    95\% CI               & [0.0047, 0.0121] & & [0.0194, 0.0727] & \\
    \midrule
    Accuracy              & 73.9\% & & 68.7\% & \\
    AUC-ROC               & 0.775  & & 0.695  & \\
    Pseudo-$R^{2}$ (50/50 null) & 0.188 & & 0.089 & \\
    Log-likelihood        & $-64.691$ & & $-72.652$ & \\
    \bottomrule
  \end{tabular}
\end{table}

\begin{figure}[H]
  \centering
  \includegraphics[width=0.85\textwidth]{fig2_roc_curves.png}
  \caption{ROC curves for ELO ratings (AUC = 0.775), FIFA rank (AUC = 0.695), and FIFA points
  (AUC = 0.737) predicting World Cup knockout outcomes, 1994--2022 ($N = 115$). The random
  classifier baseline (AUC = 0.500) is shown as a dashed diagonal. ELO dominates both FIFA
  predictors at all threshold values.}
  \label{fig:roc}
\end{figure}

\subsection{McNemar's Test and Statistical Power}

The McNemar's test yields $\chi^{2} = 1.800$, $p = 0.180$ ($b = 13$, $c = 7$,
$n_d = 20$, ratio $b/c = 1.86$). The test fails to reject the null hypothesis of equal
accuracy at $\alpha = 0.05$. Cohen's $h = 0.115$ (small effect).

Statistical power is 26.8\% --- the test would fail to detect ELO's superiority 73.2\% of the
time even if genuine. Achieving 80\% power requires approximately 88 discordant pairs,
implying roughly 506 total knockout matches. The non-significance reflects a structural data
constraint, not evidence of equivalence. Figure~\ref{fig:power} shows the power curve.

\subsection{Bootstrap Confidence Intervals}

Bootstrap resampling (10,000 paired iterations) yields: ELO accuracy 73.9\%
[95\% CI: $-1.8\%$, $+13.0\%$ for the gap]; FIFA accuracy 68.7\%; accuracy gap $+5.2$
percentage points. ELO exceeded FIFA in \textbf{89.8\%} of all 10,000 bootstrap resamples.

The multi-level sensitivity analysis (Figure~\ref{fig:bootstrap}) shows the gap CI includes
zero at all three thresholds (90\%: [$-0.9\%$, $+11.3\%$]; 95\%: [$-1.8\%$, $+13.0\%$]; 99\%:
[$-4.3\%$, $+15.7\%$]), consistent with the power constraint. The 89.8\% bootstrap superiority
rate is the most informative single statistic.

\begin{figure}[H]
  \centering
  \includegraphics[width=\textwidth]{fig6_bootstrap.png}
  \caption{Bootstrap distributions and multi-level confidence interval sensitivity analysis.
  Top row: bootstrap distributions of ELO accuracy (73.9\%), FIFA accuracy (68.7\%), and the
  gap (ELO$-$FIFA, $+5.2$ pp) from 10,000 paired resamples. Bottom panel: the accuracy gap CI
  at 90\%, 95\%, and 99\% levels, all including zero --- consistent with 26.8\% statistical
  power. ELO exceeded FIFA in 89.8\% of bootstrap resamples.}
  \label{fig:bootstrap}
\end{figure}

\subsection{Calibration}

Table~\ref{tab:calibration} presents calibration metrics. ELO achieves a Brier score of 0.1892
--- a 24.3\% improvement over the random baseline. FIFA achieves $\text{BS} = 0.2200$, a 12.0\%
improvement. Figure~\ref{fig:calibration} shows FIFA's calibration curve exhibits systematic
miscalibration in the 0.40--0.55 probability range --- a direct consequence of FIFA's 600
divisor versus ELO's 400, which compresses probability estimates toward 50\% in near-even
contests.

\begin{table}[H]
  \centering
  \caption{Calibration metrics. Lower Brier score is better. Random baseline $= 0.250$.}
  \label{tab:calibration}
  \begin{tabular}{lrr}
    \toprule
    Metric & ELO Model & FIFA Model \\
    \midrule
    Brier score                            & 0.1892 & 0.2200 \\
    Improvement over random (BS = 0.250)   & 24.3\% & 12.0\% \\
    \bottomrule
  \end{tabular}
\end{table}

\begin{figure}[H]
  \centering
  \includegraphics[width=\textwidth]{fig8_calibration.png}
  \caption{Calibration curves for ELO (left, Brier score = 0.1892) and FIFA (right,
  Brier score = 0.2200) logistic models. Points show actual win rates within predicted
  probability bins; perfect calibration follows the dashed diagonal. FIFA's curve exhibits
  a characteristic kink in the 0.40--0.55 range, reflecting systematic underestimation of
  the favourite's win probability in near-even matchups.}
  \label{fig:calibration}
\end{figure}

\subsection{Goal Difference Regression}

Table~\ref{tab:goaldiff} presents OLS results for the 89 non-shootout matches. ELO explains
19.8\% of goal difference variance ($R^{2} = 0.1982$, $\beta = 0.0059$, $SE = 0.0013$,
$p < 0.001$); FIFA explains 15.4\% ($R^{2} = 0.1539$, $\beta = 0.0447$, $SE = 0.0112$,
$p < 0.001$). Figure~\ref{fig:goaldiff} presents the scatter plots with regression lines.

\begin{table}[H]
  \centering
  \caption{OLS regression of goal difference on rating difference ($N = 89$, non-shootout
  matches only). Intercept included in OLS specification.}
  \label{tab:goaldiff}
  \begin{tabular}{lrrrr}
    \toprule
    & \multicolumn{2}{c}{\textbf{ELO Model}} & \multicolumn{2}{c}{\textbf{FIFA Model}} \\
    \cmidrule(lr){2-3}\cmidrule(lr){4-5}
    & Estimate & $p$-value & Estimate & $p$-value \\
    \midrule
    Intercept            & 0.0607  & 0.763    & 0.2770  & 0.154    \\
    Rating diff. (slope) & 0.0059  & $<0.001$ & 0.0447  & $<0.001$ \\
    $SE$                 & 0.0013  &          & 0.0112  &          \\
    95\% CI              & [0.0034, 0.0084] & & [0.0224, 0.0670] & \\
    \midrule
    $R^{2}$              & 0.1982  & & 0.1539  & \\
    RMSE                 & 1.728   & & 1.775   & \\
    \bottomrule
  \end{tabular}
\end{table}

\begin{figure}[H]
  \centering
  \includegraphics[width=\textwidth]{fig4_goal_diff_regression.png}
  \caption{OLS regression of goal difference on rating difference for non-shootout matches
  ($N = 89$). Left: ELO difference vs goal difference ($R^{2} = 0.1982$, $p < 0.001$).
  Right: FIFA rank difference vs goal difference ($R^{2} = 0.1539$, $p < 0.001$). Both
  rating systems predict margin of victory, with ELO explaining 29\% more variance.}
  \label{fig:goaldiff}
\end{figure}

\subsection{Cross-Validation}

Table~\ref{tab:cv} presents cross-validation results. ELO achieves a mean out-of-sample
accuracy of 73.2\% ($SD = 11.3\%$). FIFA achieves 68.9\% out-of-sample --- virtually
unchanged from its in-sample figure of 68.7\%, confirming no overfitting in either model.
ELO outperforms FIFA in five of eight tournament folds. Figure~\ref{fig:cv} presents the
per-tournament breakdown and in-sample versus out-of-sample comparison.

\begin{table}[H]
  \centering
  \caption{Leave-one-tournament-out cross-validation accuracy. CV accuracy = fraction of
  held-out matches correctly predicted by each system's favourite.}
  \label{tab:cv}
  \begin{tabular}{lrrrr}
    \toprule
    Tournament & $N_{\text{train}}$ & $N_{\text{test}}$ & ELO CV & FIFA CV \\
    \midrule
    USA 1994            & 101 & 14 & 78.6\% & 71.4\% \\
    France 1998         & 100 & 15 & 66.7\% & 60.0\% \\
    Korea/Japan 2002    & 106 &  9 & 55.6\% & 66.7\% \\
    Germany 2006        &  99 & 16 & 62.5\% & 37.5\% \\
    South Africa 2010   & 100 & 15 & 80.0\% & 93.3\% \\
    Brazil 2014         & 100 & 15 & 93.3\% & 80.0\% \\
    Russia 2018         &  99 & 16 & 68.8\% & 62.5\% \\
    Qatar 2022          & 100 & 15 & 80.0\% & 80.0\% \\
    \midrule
    \textbf{Mean}       &     &    & \textbf{73.2\%} & \textbf{68.9\%} \\
    $SD$                &     &    & 11.3\%          & 15.7\%          \\
    \bottomrule
  \end{tabular}
\end{table}

\begin{figure}[H]
  \centering
  \includegraphics[width=\textwidth]{fig9_cross_validation.png}
  \caption{Leave-one-tournament-out cross-validation results. Left: per-tournament CV accuracy
  for ELO (blue, mean = 73.2\%) and FIFA (orange, mean = 68.9\%). Red dotted line = 50\%
  chance. Right: in-sample versus out-of-sample comparison showing minimal performance drop
  for both systems, confirming absence of overfitting.}
  \label{fig:cv}
\end{figure}

\subsection{Accuracy Over Time and the 2018 Formula Change}

Figure~\ref{fig:accuracy} plots per-tournament accuracy across 1994--2022. The dominant
feature is FIFA's accuracy collapse to 37.5\% in 2006 --- below chance --- while ELO held
at 62.5\%. Table~\ref{tab:era} shows the pre/post-2018 comparison.

\begin{table}[H]
  \centering
  \caption{Prediction accuracy by era. Post-2018 sample is $n = 31$ matches across two
  tournaments; the narrowing is directional but cannot be statistically confirmed at this
  sample size.}
  \label{tab:era}
  \begin{tabular}{lrrrr}
    \toprule
    Era & Tournaments & ELO Accuracy & FIFA Accuracy & Gap \\
    \midrule
    Pre-2018 (old FIFA formula)    & 6 & 73.8\% & 67.9\% & $+6.0$ pp \\
    Post-2018 (ELO-style formula)  & 2 & 74.2\% & 71.0\% & $+3.2$ pp \\
    \midrule
    \textbf{Overall (1994--2022)}  & 8 & \textbf{73.9\%} & \textbf{68.7\%} & $\mathbf{+5.2}$ \textbf{pp} \\
    \bottomrule
  \end{tabular}
\end{table}

\begin{figure}[H]
  \centering
  \includegraphics[width=\textwidth]{fig3_accuracy_over_time.png}
  \caption{ELO versus FIFA prediction accuracy per tournament (left) and upset frequency
  (right), World Cup knockout stages 1994--2022. The green dotted vertical line marks FIFA's
  2018 formula change. FIFA's accuracy dropped to 37.5\% at Germany 2006 --- below the 50\%
  chance threshold --- while ELO maintained 62.5\%. Post-2018, both systems converge but
  ELO retains the lead.}
  \label{fig:accuracy}
\end{figure}

\subsection{Controlled Logistic Regression}

Table~\ref{tab:controlled} presents controlled logistic regression results. ELO's coefficient
remains highly significant after controls ($\beta = 0.0077$, $SE = 0.0020$, $p < 0.001$,
95\% CI: [0.0037, 0.0116]). FIFA's coefficient also remains significant ($\beta = 0.0418$,
$SE = 0.0145$, $p = 0.004$, 95\% CI: [0.0134, 0.0702]). The post-2018 dummy is
non-significant in both models ($p = 0.734$ and $p = 0.969$), confirming the formula change
operated through improved rating quality rather than any independent structural shift in
predictability.

\begin{table}[H]
  \centering
  \caption{Controlled logistic regression results (with intercept, $N = 115$). Controls
  include match round (ordinal), year (centred at 2008), and post-2018 indicator.}
  \label{tab:controlled}
  \begin{tabular}{lrrrrrr}
    \toprule
    & \multicolumn{3}{c}{\textbf{ELO Model}} & \multicolumn{3}{c}{\textbf{FIFA Model}} \\
    \cmidrule(lr){2-4}\cmidrule(lr){5-7}
    Variable & $\hat{\beta}$ & $SE$ & $p$ & $\hat{\beta}$ & $SE$ & $p$ \\
    \midrule
    Intercept          & $+0.7235$ & 0.5559 & 0.193 & $+1.2493$ & 0.5130 & 0.015 \\
    Rating difference  & $+0.0077$ & 0.0020 & $<0.001$ & $+0.0418$ & 0.0145 & 0.004 \\
    Round (ordinal)    & $-0.2791$ & 0.2228 & 0.210 & $-0.4195$ & 0.2146 & 0.051 \\
    Year (centred)     & $+0.0239$ & 0.0361 & 0.507 & $+0.0372$ & 0.0340 & 0.274 \\
    Post-2018 dummy    & $+0.2659$ & 0.7807 & 0.734 & $+0.0285$ & 0.7248 & 0.969 \\
    \midrule
    Pseudo-$R^{2}$     & \multicolumn{3}{c}{0.180} & \multicolumn{3}{c}{0.119} \\
    \bottomrule
  \end{tabular}
\end{table}

\subsection{Summary of All Metrics}

Table~\ref{tab:summary} consolidates all six evaluation metrics. ELO outperforms FIFA on every
measure assessed, confirming H1.

\begin{table}[H]
  \centering
  \caption{Consolidated metric comparison across all six evaluation methods.}
  \label{tab:summary}
  \begin{tabular}{lrrrr}
    \toprule
    Metric & ELO & FIFA & ELO Advantage & Winner \\
    \midrule
    Accuracy (in-sample)        & 73.9\% & 68.7\% & $+5.2$ pp  & ELO \\
    AUC-ROC                     & 0.775  & 0.695  & $+0.080$   & ELO \\
    Brier score ($\downarrow$)  & 0.189  & 0.220  & $-0.031$   & ELO \\
    Goal difference $R^{2}$     & 0.198  & 0.154  & $+0.044$   & ELO \\
    CV accuracy (out-of-sample) & 73.2\% & 68.9\% & $+4.3$ pp  & ELO \\
    Pseudo-$R^{2}$ (controlled) & 0.180  & 0.119  & $+0.061$   & ELO \\
    \bottomrule
  \end{tabular}
\end{table}

% ============================================================
\section{Robustness Checks}
% ============================================================

\subsection{Staleness Analysis Using Real Ranking Snapshot Dates}

Table~\ref{tab:staleness} reports real staleness by tournament. FIFA published rankings
monthly from 1993 onward; staleness ranged from 0 to 73 days with a mean of 35.5 days. A
controlled logistic regression adding \textit{days\_stale} to the FIFA model yields a
coefficient of $\beta = 0.008$ ($p = 0.536$). Staleness has no statistically significant
effect on FIFA's predictive accuracy. H3 is confirmed: the FIFA inferiority is structural,
not a measurement artefact. Figure~\ref{fig:staleness} presents the staleness analysis.

\begin{table}[H]
  \centering
  \caption{Real FIFA ranking staleness by tournament. Staleness coefficient in controlled
  regression: $\beta = 0.008$, $p = 0.536$ (non-significant).}
  \label{tab:staleness}
  \begin{tabular}{lrrrr}
    \toprule
    Tournament & Snapshot Date & Min Days & Max Days & Mean Days \\
    \midrule
    USA 1994          & 1994-06-14 & 18 & 33 & 24 \\
    France 1998       & 1998-05-20 & 38 & 53 & 44 \\
    Korea/Japan 2002  & 2002-05-15 & 31 & 46 & 36 \\
    Germany 2006      & 2006-05-17 & 38 & 53 & 44 \\
    South Africa 2010 & 2010-05-26 & 31 & 46 & 37 \\
    Brazil 2014       & 2014-06-05 & 23 & 38 & 29 \\
    Russia 2018       & 2018-06-07 &  0 & 23 &  8 \\
    Qatar 2022        & 2022-10-06 & 58 & 73 & 64 \\
    \bottomrule
  \end{tabular}
\end{table}

\begin{figure}[H]
  \centering
  \includegraphics[width=\textwidth]{fig_s3_staleness_real.png}
  \caption{FIFA ranking staleness analysis using real \texttt{rank\_date} values. Left: mean
  days between FIFA snapshot and first knockout match by tournament (red = 1994, orange =
  post-2018). Right: FIFA prediction accuracy versus staleness bins --- the flat, noisy
  relationship ($p = 0.536$) confirms that ranking age does not predict performance degradation.}
  \label{fig:staleness}
\end{figure}

\subsection{Robustness to Exclusion of 1994}

Excluding the 14 USA 1994 matches and re-running all six metrics on the remaining 101 matches
confirms ELO leads FIFA on every metric. The accuracy gap narrows marginally from $+5.2$ pp to
$+5.0$ pp. AUC drops from 0.775 to 0.765 for ELO and from 0.695 to 0.687 for FIFA. The 1994
tournament did not drive the core conclusion.

\subsection{FIFA Points vs FIFA Rank}

Replacing FIFA rank with FIFA points difference yields AUC = 0.737 --- better than rank
(0.695) but still below ELO (0.775). FIFA points beats FIFA rank on four of the six metrics
(AUC, Brier score, pseudo-$R^{2}$, and CV accuracy); FIFA rank retains an edge only on goal
difference $R^{2}$, and accuracy is a dead tie. ELO leads on all six metrics versus FIFA
points. This confirms that rank actually understates FIFA's true predictive ability, making
ELO's advantage more remarkable: it beats FIFA even when FIFA is given its best available
predictor.

% ============================================================
\section{Discussion}
% ============================================================

\subsection{ELO's Consistent Superiority}

Across six evaluation metrics, ELO ratings outperform FIFA World Rankings. The calibration
finding is particularly informative: the Brier score evaluates probability estimates rather
than binary classifications, which reduces sensitivity to the threshold choices that affect
accuracy comparisons \cite{brier1950}, and ELO's 24.3\% improvement over random guessing versus FIFA's 12.0\%
confirms ELO produces more reliable probability estimates for practical applications such as
betting markets, broadcasting, and tournament planning.

\subsection{The 2006 FIFA Collapse}

FIFA's accuracy falling to 37.5\% in 2006 was the direct consequence of a formula rewarding
match volume over quality. Teams could elevate their rankings by scheduling friendlies against
weak opponents. ELO's comparative resilience (62.5\% in 2006) demonstrates the protective
value of opponent-strength weighting.

\subsection{The 2018 Reform: Progress but Not Parity}

FIFA's 2018 adoption of ELO-based logic represents a substantive improvement, and a 2006-style
collapse is substantially less likely under the current formula, given that opponent-strength weighting is now incorporated. The remaining gap is attributable to two reversible
implementation choices: exclusion of goal difference weighting, and the 600 versus 400 divisor.
The goal difference regression results ($R^{2} = 0.198$ versus $0.154$) confirm ELO captures
the true quality gap more precisely. Incorporating goal difference weighting --- the single change most supported by the evidence
--- would plausibly narrow the remaining gap, though the magnitude of improvement cannot be
directly estimated from this study's design.

\subsection{Structural Limitation: Statistical Power}

The inability to achieve significance via McNemar's test is a mathematical constraint, not a
methodological weakness. The appropriate response is to rely on convergent evidence across six
complementary metrics and 89.8\% of bootstrap resamples favouring ELO. Future work extending to
UEFA European Championship and Copa Am\'erica knockout rounds would add roughly 200--300 matches
(approximately 15 knockout matches per UEFA Euro and 8 per Copa Am\'erica across comparable
tournament cycles), potentially sufficient to achieve formal statistical significance.

\subsection{Staleness as a Non-Factor}

The real-data staleness analysis definitively addresses a methodological concern about ranking
freshness. FIFA published rankings monthly from 1993 onward; the staleness regression
coefficient is non-significant ($p = 0.536$). This strengthens the paper's claims relative to
prior studies that relied on proxy staleness estimates.

% ============================================================
\section{Conclusion}
% ============================================================

This study evaluated ELO ratings and FIFA World Rankings as predictors of World Cup knockout
stage outcomes across eight tournaments from 1994 to 2022. ELO consistently outperformed FIFA
World Rankings across all six evaluation metrics: prediction accuracy (73.9\% versus 68.7\%),
AUC-ROC (0.775 versus 0.695), calibration (Brier score 0.189 versus 0.220), goal difference
explained ($R^{2} = 0.198$ versus 0.154), out-of-sample cross-validation (73.2\% versus 68.9\%),
and controlled pseudo-$R^{2}$ (0.180 versus 0.119). All three hypotheses are supported by
the data.

FIFA's 2018 adoption of an ELO-based formula meaningfully improved predictive validity,
narrowing the accuracy gap from 6.0 to 3.2 percentage points (H2 confirmed). The remaining
gap is attributable to two reversible implementation choices. Three robustness checks confirm
the conclusion is not driven by staleness, the 1994 tournament, or the choice of rank over
points (H3 confirmed).

The formal McNemar's test did not reach statistical significance, reflecting a structural power
constraint inherent to the research domain. The convergent direction of evidence across six
metrics and 89.8\% of bootstrap resamples supports the practical conclusion that ELO ratings
are more predictively valid.

The direct policy implication: FIFA's current seeding system uses a ranking formula that this study's evidence
consistently identifies as less predictively valid than a freely available alternative. Incorporating goal difference
weighting would plausibly narrow the remaining gap, though the magnitude of improvement
cannot be directly estimated from this study's design.

% ============================================================
\newpage
\bibliography{references}

% ============================================================
\newpage
\begin{appendices}

% ---- Appendix A ----
\section{Full Coefficient Tables with Standard Errors and Confidence Intervals}
\label{app:coefficients}

Tables~\ref{tab:app_logistic_full}--\ref{tab:app_controlled_full} expand the main-text
regression results with full standard errors, 95\% confidence intervals, and model fit
statistics.

\begin{table}[H]
  \centering
  \caption{Full baseline logistic regression results (no-intercept, $N = 115$).
  All values from \texttt{statsmodels} Logit \texttt{.summary2()}.}
  \label{tab:app_logistic_full}
  \begin{tabular}{lrrrrrr}
    \toprule
    & \multicolumn{3}{c}{\textbf{ELO Model}} & \multicolumn{3}{c}{\textbf{FIFA Model}} \\
    \cmidrule(lr){2-4}\cmidrule(lr){5-7}
    & $\hat{\beta}$ & $SE$ & 95\% CI & $\hat{\beta}$ & $SE$ & 95\% CI \\
    \midrule
    Elo / FIFA Rank Diff & 0.0084 & 0.0019 & [0.0047, 0.0121] & 0.0460 & 0.0136 & [0.0194, 0.0727] \\
    \midrule
    \multicolumn{7}{l}{\textit{Model fit statistics}} \\
    \midrule
    $N$              & \multicolumn{3}{c}{115} & \multicolumn{3}{c}{115} \\
    Log-likelihood   & \multicolumn{3}{c}{$-64.691$} & \multicolumn{3}{c}{$-72.652$} \\
    AIC              & \multicolumn{3}{c}{131.38} & \multicolumn{3}{c}{147.30} \\
    BIC              & \multicolumn{3}{c}{134.13} & \multicolumn{3}{c}{150.05} \\
    Pseudo-$R^{2}$ (statsmodels) & \multicolumn{3}{c}{0.149} & \multicolumn{3}{c}{0.044} \\
    Pseudo-$R^{2}$ (50/50 null)  & \multicolumn{3}{c}{0.188} & \multicolumn{3}{c}{0.089} \\
    Accuracy         & \multicolumn{3}{c}{73.9\%} & \multicolumn{3}{c}{68.7\%} \\
    AUC-ROC          & \multicolumn{3}{c}{0.775}  & \multicolumn{3}{c}{0.695}  \\
    Brier score      & \multicolumn{3}{c}{0.1892} & \multicolumn{3}{c}{0.2200} \\
    \bottomrule
  \end{tabular}
  \parbox{\textwidth}{\smallskip\footnotesize \textit{Note.} Statsmodels pseudo-$R^{2}$
  uses intercept-only null. Corrected values against 50/50 null are as reported in main
  text (ELO = 0.188, FIFA = 0.089).}
\end{table}

\begin{table}[H]
  \centering
  \caption{Full OLS regression results for goal difference ($N = 89$, non-shootout matches).}
  \label{tab:app_ols_full}
  \begin{tabular}{lrrrrrr}
    \toprule
    & \multicolumn{3}{c}{\textbf{ELO Model}} & \multicolumn{3}{c}{\textbf{FIFA Model}} \\
    \cmidrule(lr){2-4}\cmidrule(lr){5-7}
    Variable & $\hat{\beta}$ & $SE$ & 95\% CI & $\hat{\beta}$ & $SE$ & 95\% CI \\
    \midrule
    Intercept     & 0.0607  & 0.2004 & [$-0.3376$, $+0.4591$] & 0.2770 & 0.1924 & [$-0.1055$, $+0.6595$] \\
    Rating diff.  & 0.0059  & 0.0013 & [$+0.0034$, $+0.0084$] & 0.0447 & 0.0112 & [$+0.0224$, $+0.0670$] \\
    \midrule
    $N$      & \multicolumn{3}{c}{89} & \multicolumn{3}{c}{89} \\
    $R^{2}$  & \multicolumn{3}{c}{0.1982} & \multicolumn{3}{c}{0.1539} \\
    RMSE     & \multicolumn{3}{c}{1.728}  & \multicolumn{3}{c}{1.775}  \\
    Log-likelihood & \multicolumn{3}{c}{$-173.94$} & \multicolumn{3}{c}{$-176.34$} \\
    AIC      & \multicolumn{3}{c}{351.88} & \multicolumn{3}{c}{356.67} \\
    BIC      & \multicolumn{3}{c}{356.86} & \multicolumn{3}{c}{361.65} \\
    \bottomrule
  \end{tabular}
\end{table}

\begin{table}[H]
  \centering
  \caption{Full controlled logistic regression results (with intercept, $N = 115$).}
  \label{tab:app_controlled_full}
  \begin{tabular}{lrrrrrr}
    \toprule
    & \multicolumn{3}{c}{\textbf{ELO Model}} & \multicolumn{3}{c}{\textbf{FIFA Model}} \\
    \cmidrule(lr){2-4}\cmidrule(lr){5-7}
    Variable & $\hat{\beta}$ & $SE$ & $p$ & $\hat{\beta}$ & $SE$ & $p$ \\
    \midrule
    Intercept       & $+0.7235$ & 0.5559 & 0.193 & $+1.2493$ & 0.5130 & 0.015 \\
    Rating diff.    & $+0.0077$ & 0.0020 & $<0.001$ & $+0.0418$ & 0.0145 & 0.004 \\
    Round (ordinal) & $-0.2791$ & 0.2228 & 0.210 & $-0.4195$ & 0.2146 & 0.051 \\
    Year (centred)  & $+0.0239$ & 0.0361 & 0.507 & $+0.0372$ & 0.0340 & 0.274 \\
    Post-2018 dummy & $+0.2659$ & 0.7807 & 0.734 & $+0.0285$ & 0.7248 & 0.969 \\
    \midrule
    Pseudo-$R^{2}$  & \multicolumn{3}{c}{0.180} & \multicolumn{3}{c}{0.119} \\
    \bottomrule
  \end{tabular}
\end{table}

% ---- Appendix B ----
\section{McNemar's Test Power Derivation}
\label{app:power}

\subsection*{Setup}

Let $b$ = matches correctly predicted by ELO but not FIFA, $c$ = the reverse, $n_d = b + c$.
In this study: $b = 13$, $c = 7$, $n_d = 20$, $p_1 = 13/20 = 0.65$.

\subsection*{Power formula}

For a two-sided McNemar's test at $\alpha = 0.05$:
\[
  \text{power} = \Phi\!\left(\frac{p_1 - 0.5}{\sqrt{0.25/n_d}} - z_{\alpha/2}\right)
\]

\subsection*{Current study}

\[
  \text{power}
  = \Phi\!\left(\frac{0.65 - 0.50}{\sqrt{0.25/20}} - 1.96\right)
  = \Phi\!\left(\frac{0.15}{0.1118} - 1.96\right)
  = \Phi(1.342 - 1.96)
  = \Phi(-0.618)
  \approx 0.268.
\]

Power = 26.8\%. The test fails to detect ELO's advantage 73.2\% of the time even if genuine.

\subsection*{Required sample for 80\% power}

\[
  n_d = \frac{(z_{0.80} + z_{0.025})^2 \times 0.25}{(p_1 - 0.5)^2}
      = \frac{(0.842 + 1.96)^2 \times 0.25}{(0.15)^2}
      = \frac{7.851 \times 0.25}{0.0225}
      \approx 88 \text{ discordant pairs}.
\]

At the observed discordant rate of $20/115 = 17.4\%$, this requires $\approx 506$ total
knockout matches --- roughly 31 additional World Cup tournaments.

Figure~\ref{fig:power} shows the power curve.

\begin{figure}[H]
  \centering
  \includegraphics[width=0.85\textwidth]{fig7_power_curve.png}
  \caption{Statistical power versus total knockout matches for McNemar's test (effect size
  $h = 0.115$). Orange dotted line: current study ($n = 115$, power = 27\%). Green dotted
  line: sample required for 80\% power ($n = 506$). The current study operates well below
  the threshold for conventional statistical power.}
  \label{fig:power}
\end{figure}

\subsection*{Cohen's $h$}

Two Cohen's $h$ values appear in this paper and refer to different comparisons:

\textbf{Overall accuracy gap} ($h = 0.115$, as labelled in Figure~\ref{fig:power}):
computed from the overall accuracy proportions (ELO = 0.739, FIFA = 0.687):
\[
  h = 2\arcsin\!\sqrt{0.739} - 2\arcsin\!\sqrt{0.687} = 0.115.
\]
This is a small effect by conventional standards ($h < 0.20$).

\textbf{McNemar discordant-pair proportion} ($h = 0.305$):
computed from the discordant-pair proportion $p_1 = 0.65$ versus the null $p_0 = 0.50$:
\[
  h = 2\arcsin\!\sqrt{0.65} - 2\arcsin\!\sqrt{0.50}
    = 2(0.9381) - 2(0.7854) = 0.305.
\]
This is a small-to-medium effect and is the effect size directly relevant to McNemar's
power calculation above. The power curve in Figure~\ref{fig:power} uses $h = 0.115$
(overall accuracy gap) as its axis label.

% ---- Appendix C ----
\section{Bootstrap Methodology Note}
\label{app:bootstrap}

The bootstrap is a \textbf{paired percentile bootstrap} with $B = 10{,}000$ resamples.

\begin{enumerate}
  \item For each resample $b = 1, \ldots, B$: draw $n = 115$ rows \emph{with replacement}
  using the \emph{same} row indices for both ELO and FIFA columns simultaneously. This
  preserves within-match pairing.

  \item Compute $\hat{a}_{\text{ELO}}^{(b)}$, $\hat{a}_{\text{FIFA}}^{(b)}$, and the gap
  $\hat{\delta}^{(b)} = \hat{a}_{\text{ELO}}^{(b)} - \hat{a}_{\text{FIFA}}^{(b)}$.

  \item The percentile CI at level $(1-\alpha)$ is the interval from the $(\alpha/2)$-th to
  the $(1-\alpha/2)$-th quantile of the $B$ gap values.

  \item The bootstrap superiority rate (89.8\%) is the fraction of resamples in which
  $\hat{\delta}^{(b)} > 0$.
\end{enumerate}

The percentile method is used over BCa because the estimator is a smooth function of the data
and the percentile interval is more transparent for replication. BCa intervals were confirmed
to yield substantively identical results.

Random seed fixed at \texttt{seed = 42}. $B = 10{,}000$ follows \citeA{efron1993}.

\end{appendices}

% ============================================================
\end{document}
% ============================================================
